% Standalone insertion; requires amsmath, amsthm and a proposition environment.
\section{An exact finite-width tanh certificate}
\label{sec:finite}

The following calculation identifies a concrete obstruction to extending a
width-one online-mistake proof merely by replicating its neurons. It does
not establish the prediction of the unrestricted wider-network posterior:
a different wider implementation of a wrong function could still have
smaller energy.

Use $d=7$ input bits $(x_0,x_1,x_2,y_0,y_1,y_2,z)$, seven tanh hidden
layers of common width $N$, and a linear readout. There is one appended
constant input and no other biases. The Gaussian quadratic energy is
\[
 \kappa_N=\frac7N\|W^1\|_F^2+
       \sum_{\ell=2}^{7}\|W^\ell\|_F^2+N\|v\|_2^2.
\]
Put $X=\sum_{i=0}^2 2^i(x_i+1)/2$, $Y=\sum_{i=0}^2 2^i(y_i+1)/2$ and
\[
 f(x,y,z)=+1\ \Longleftrightarrow\ z=-1\text{ and }X\ge Y.
\]
After the calibration in Section~\ref{sec:selector} and queries
$X=3,4,5,6$, $Y=0$, $z=+1$, the affine minimum is the sign function
\[
 g_6=\operatorname{sign}\!\left(-6-7z+
                       \sum_{i=0}^2 2^i(x_i-y_i)\right).
\]
Its minimum squared affine coefficient norm at unit global margin is
$Q(g_6)=42+36+49=127$. In particular, it wrongly labels the next query
$(X,Y,z)=(7,0,+1)$ positive.

\begin{proposition}
For every integer $N\ge64$, the target $f$ has a width-$N$ implementation
with all weights of absolute value less than three, global output margin
greater than $1/2$, and
\[
 \kappa_N<15.7.
\]
Every width-one chain computing $g_6$ with global output margin at least
$1/2$ has energy at least
\[
 8(7\cdot127/4)^{1/8}>15.7.
\]
Equal replication of that scalar chain to width $N$ preserves its energy.
Thus no equally replicated scalar minimizer for this history minimizes
energy among all width-$N$ networks.
\end{proposition}
\begin{proof}
Let $n=\lfloor N/64+1/2\rfloor$ and $p=n/N$. For $N\ge64$,
\[
 1/95\le p\le1/48.
\]
Indeed, when $n=1$ we have $64\le N\le95$; when $n\ge2$, the inequalities
$64n-32\le N\le64n+31$ imply the displayed bounds.

Write
\[
 U=1+\sum_{i=0}^2 2^i(x_i-y_i),\quad
 \alpha=\frac{55}{100},\quad b=\frac{17}{100},\quad
 r=\frac{147}{1000},\quad s=\frac{675}{1000}.
\]
The first $n$ neurons of layer one compute $a=\tanh(\alpha U)$;
the remaining $N-n$ compute $c=\tanh(b(z+1))$. Every neuron of layer two
computes the same value $h_2=\tanh(ra-sc)$, using weights $r/n$ and
$-s/(N-n)$ on the two groups. The next five layers are equal replicas
with respective effective scalar weights
\[
 (w_3,w_4,w_5,w_6,w_7)
   =\frac1{1000}(1361,1365,1373,1389,1419).
\]
Each physical weight in layer $j$ is $w_j/N$. The readout has $N$ equal
weights $(3/2)/N$.

When $z=-1$, $c=0$ and $|U|\ge1$, so the second-layer preactivation
has the sign of $U$ and absolute value at least $r\tanh\alpha$.
When $z=+1$, it is at most $r-s\tanh(2b)<0$.
Consequently the network has the target's sign everywhere, and the
minimum second-layer preactivation magnitude is at least
\[
 m=\min\{r P(\alpha),\ sP(2b)-r\},\qquad P(u)=u-u^3/3.
\]
Here $\tanh u\ge P(u)$ for $u\ge0$: the difference vanishes at zero
and has derivative $u^2-\tanh^2u\ge0$.
Starting with $\underline h_2=P(m)$ and recursively setting
$\underline h_j=P(w_j\underline h_{j-1})$ gives positive rational
lower bounds on the hidden magnitudes. Exact rational arithmetic gives
\[
 \frac32\underline h_7>\frac12
 \qquad\text{(numerically, }0.5033684798\ldots\text{)}.
\]
The accompanying script \texttt{verification/check\_finite\_certificate.py}
verifies this inequality without floating-point arithmetic.

Since the squared affine coefficient norm of $U$ is $43$, the energy is
exactly
\[
 E(p)=7\bigl(43p\alpha^2+2(1-p)b^2\bigr)
       +\frac{r^2}{p}+\frac{s^2}{1-p}
       +\sum_{j=3}^7w_j^2+\frac94.
\]
This is convex in $p$. Exact rational evaluation at the two endpoints
shows
\[
 E(1/95)<15.666,\qquad E(1/48)<15.590,
\]
so $E(p)<15.7$ on the entire interval. The largest physical weight in
this implementation is $4\alpha=2.2$.

For the scalar lower bound, let $\rho$ be its least absolute first-layer
preactivation. Its first coefficient vector has squared norm at least
$127\rho^2$. Since $|\tanh u|\le|u|$, the product of $\rho$ and the
absolute values of its remaining seven weights is at least $1/2$.
Applying arithmetic--geometric mean to the eight energy terms yields
$\kappa_1\ge8(7\cdot127/4)^{1/8}$. The strict comparison with $15.7$
is the rational inequality $(157/80)^8<889/4$.
\end{proof}
